% Source: https://artofproblemsolving.com/community/c3h3129147_2022_midmichigan_mathematical_olympiad__grades_56
% Stable identity: aops-topic-3129147-q03
Let the original vertices be $V_0,\ldots,V_5$ in order, with centre $O$. Start with the three congruent quadrilateral sectors $OV_0V_1V_2$, $OV_2V_3V_4$, and $OV_4V_5V_0$. Replace each radial boundary $OV_{2j}$ by a slightly bent two-segment path, choosing the three bends as rotations of one another by $120^\circ$. For sufficiently small nonzero bends the paths meet only at $O$ and remain inside the hexagon. Each sector now has six corners: $O$, the first bend, two consecutive outer vertices, the third outer vertex, and the second bend. Rotations by $120^\circ$ carry each region to the next, proving congruence. The explicit construction below makes all six corners noncollinear with their neighbours.
\begin{asy}
usepackage("amsmath");
size(7cm);
pair O=(0,0);pair[] V,K;for(int i=0;i<6;++i)V.push(dir(60*i));for(int j=0;j<3;++j)K.push(0.5dir(120*j)+0.08dir(120*j+90));filldraw(O--K[0]--V[0]--V[1]--V[2]--K[1]--cycle,lightblue,black);filldraw(O--K[1]--V[2]--V[3]--V[4]--K[2]--cycle,lightgreen,black);filldraw(O--K[2]--V[4]--V[5]--V[0]--K[0]--cycle,lightyellow,black);dot(O);
\end{asy}
